An einem sonnigen Tag wird der Sand am Strand so heiss, dass man kaum barfuss darüber gehen kann, während das Meerwasser angenehm kühl bleibt. Nachts kühlt der Sand rasch ab, das Wasser kaum.

Sand und Wasser erhalten pro Quadratmeter etwa gleich viel Sonnenstrahlung. Nach $Q = c\,m\,\Delta\vartheta$ ändert sich bei gleicher Wärme die Temperatur desjenigen Stoffes stärker, der die kleinere spezifische Wärme hat. Die spezifische Wärme von Sand ist also $\result{\text{kleiner}}$ als diejenige von Wasser (trockener Sand etwa \cSaO, Wasser \cWO).

Der Unterschied wird noch verstärkt, weil das Wasser die Wärme durch Strömungen auf eine dicke Schicht verteilt und durch Verdunsten gekühlt wird, während sich im schlecht leitenden Sand nur die oberste Schicht erwärmt.
